\documentclass[12pt]{article}
\usepackage{graphicx}
\usepackage{amsmath,amsfonts,amssymb}
\usepackage[left=3cm,right=3cm,top=2cm,bottom=2cm]{geometry}
\usepackage{booktabs}
\usepackage{float}
\usepackage{hyperref}

\title{An Extended Report on the Design of a Course on Signal Processing and Machine Learning}
\author{Manas Patil \\Towards Seminar under the guidance of Prof. Vikram Gadre}
\date{05/01/2026}

\begin{document}
\maketitle

\section{Introduction}
In the previous report, we confirmed the necessity of a common language and basis of Linear Algebra to bridge concepts in Signal Processing and Machine Learning. Furthermore, we discussed a loose design of 14 hours into three sections and a brief idea on how they can be resolved to higher durations in a telescopic manner. We also had some preliminary ideas as to what exercises each section shall contain. \\
In a bid to better understand the specifics of what the lectures will contain and how they will realize this connection, this report aims to give a draft of the specific topics the series of lectures will cover, in the 1 month (4 week) duration. In that framework of this report, thus, a significant consideration is given to thinking about the exercises to be put in the course, their goals, format and manner, in which they will appear in the course. 

\section{Exercises of the Course}
Part of what makes a course from an elite institution and instructor really shine and stick through is its choice of exercises. The consistent focus on solving challenging problems right in the class to tutorial problems at students' leisure in EE603, and even other classes that I took, makes me ruminate on this section a lot. Having solved the absolute minimum of "What is the equation of this linear regressor?", "Answer in brief" and "Explain what is X." kind of questions before, the quality, tone and tenor of questions at IITB made me pleasantly surprised. I was amazed at how much I could learn from solving just 10 good questions than drilling through 30 cookie-cutter ones. Of course, the logistics of a NPTEL course prevent us to liberally challenge the learner as we wish, more so accounting for the variety of learners the course can attract. Thus the following goals of an exercise section on each of the modular stages of the course emerge:
\begin{enumerate}
    \item 
\end{enumerate}

\end{document}